\documentclass[12pt]{article}
\usepackage[utf8]{inputenc}
\usepackage[T1]{fontenc}
\usepackage[margin=1in]{geometry}
\usepackage{setspace}
\usepackage{xcolor}
\usepackage[colorlinks=true, urlcolor=blue!60!black, linkcolor=blue!60!black]{hyperref}
\setlength{\parindent}{0em}
\setlength{\parskip}{0.6em}
\renewcommand{\baselinestretch}{1.0}
\begin{document}
\begin{center}
{\LARGE \textit{Whole Woman's Health v. Jackson}, 595 U.S. 30 (2021)\par}
\vspace{0.5em}
{\large Case Brief --- Two-Page Edition\par}
\end{center}

\section{Memory Jogger}
Sovereign immunity and traditional equity barred federal suits against state judges, clerks, and the Texas Attorney General to block S.B. 8's private-enforcement abortion ban; only claims against licensing officials survived (\href{https://www.supremecourt.gov/opinions/21pdf/21-463_3ebh.pdf}{Slip Op.}).

\section{Detailed Case Facts}
Texas's S.B. 8 banned abortion after about six weeks, barred state officials from enforcing it, and let any private person sue abortion providers or helpers for at least \$10,000, with procedural rules stacked against defendants. Providers sued a state judge, a clerk, the Attorney General, licensing officials, and a private activist.

\section{Procedural History}
The district court denied dismissal; the Fifth Circuit stayed proceedings; the Supreme Court denied emergency relief 5--4, then granted certiorari before judgment.

\section{Judicial Votes}
Gorsuch for the Court (Thomas, Alito, Kavanaugh, Barrett): 9--0 dismissing the judge and activist, 5--4 dismissing the clerk and Attorney General, 8--1 (Thomas dissenting) keeping licensing officials. Roberts and Sotomayor, each joined by others, dissented in part.

\section{Holding}
\textit{Ex parte Young} permits suits only against executive officials with authority to enforce the law. Courts may not enjoin state judges, clerks, officials lacking enforcement power, unnamed private persons, or ``laws themselves.''

\section{Analysis of Opinions}
\textbf{Gorsuch:} Judges are not adverse to litigants, and enjoining a state court's machinery offends dual sovereignty; clerks are ministerial; federal defenses remain available in state court. \textbf{Thomas:} S.B. 8 stripped even licensing officials of enforcement power. \textbf{Roberts:} clerks and Paxton were proper defendants, and S.B. 8 threatens the Court's role. \textbf{Sotomayor:} the ruling invites other States to copy S.B. 8's nullification scheme.

\section{Examples: Future Applications}
\textbf{Same-Side:} Suits against judges and clerks to block a private gun-bounty law must be dismissed. \textbf{Opposite-Side:} A commissioner required to revoke accreditation of violators can be sued. \textbf{Fence-Sitter:} A registrar who must certify compliance before bounty suits may be executive or merely ministerial.

\section{Critique}
Formalists praise its fidelity to traditional equity and Article III. Fallon and others say it cheapened constitutional remedies and spurred ``vigilante federalism.''

\section{Key Quotations}
\begin{quote}
``Consistent with historical practice, a federal court exercising its equitable authority may enjoin named defendants from taking specified unlawful actions. But under traditional equitable principles, no court may `lawfully enjoin the world at large,' \dots{} or purport to enjoin challenged `laws themselves.' ''
\end{quote}
\noindent --- Justice Gorsuch, Opinion of the Court.

\end{document}
